\section{Infrastructure and Institutional Support}

The research of advanced functional nano materials relies critically on the availability of cutting-edge infrastructure and strong support capable of integrating synthesis, processing, and multiscale characterization. In Brazil, the Centro Nacional de Pesquisa em Energia e Materiais (CNPEM) provides a unique and comprehensive scientific environment that fulfills these requirements. CNPEM hosts state-of-the-art laboratories dedicated to materials research and development, most notably the Laboratório Nacional de Nanotecnologia (LNNano) and the Laboratório Nacional de Luz Síncrotron (LNLS), which operates the fourth-generation synchrotron light source Sirius. Sirius is internationally recognized as one of the most advanced synchrotron facilities in the world, offering ultra-high brilliance, coherence, and stability. The integration of the LNNano materials synthesis and processing capabilities with the advanced characterization techniques available at LNLS enables a closed and highly efficient research cycle, in which new materials can be designed, fabricated, and analyzed from the atomic to the macroscopic scale.

Within this framework, the Nanoceramics Processing Laboratory at LNNano plays a central role in the preparation and optimization of ceramic SOFC/SOEC electrolytes. The laboratory is specifically designed to process advanced ceramic systems for applications in energy conversion, catalysis, and environmental technologies, among other fields. Its infrastructure encompasses a complete set of equipment for both materials processing and physicochemical characterization. Processing capabilities include a high-energy mill for particle size reduction, which enhances reactivity and facilitates solid-state reactions. High-temperature ovens capable of operating above 1500$^\circ$C enable calcination and sintering processes essential for phase stabilization and densification. A hot press equipped with different molds allows the fabrication of dense ceramic bodies under simultaneous pressure and temperature, minimizing porosity and improving mechanical integrity. Complementary infrastructure includes precision scales for density determination and electrical characterization systems based on impedance spectroscopy, which provide insights into ionic and electronic transport mechanisms and grain boundary contributions.

At the nanoscale, the Transmission Electron Microscopy Laboratory at LNNano provides structural and spectroscopic characterization with atomic resolution. The facility is equipped with advanced transmission electron microscopes featuring high-sensitivity cameras for image acquisition and electron diffraction analysis, as well as spectrometers for detecting characteristic X-ray photons emitted by the sample. These capabilities enable the precise determination of particle size, morphology, crystalline structure, and elemental distribution.

Complementing these capabilities, the QUATI experimental facility at LNLS provides advanced time and space-resolved spectroscopic measurements under operational conditions. QUATI allows the coupling of X-ray absorption spectroscopy with Raman spectroscopy and electrochemical impedance spectroscopy, enabling simultaneous access to structural, vibrational, and electrochemical information. This approach is valuable for addressing complex characterization challenges associated with nanomaterials, especially when investigated under realistic working environments. The ability to monitor structural and electronic evolution in-situ and operando makes QUATI ideally suited for the nanomaterials studied in the present project, ensuring a comprehensive evaluation of their functional performance.
